\section{Conclusion}
\label{sec:conclusion}

We investigated how compound concepts like ``washing machine'' are represented in the residual stream of \gpttwo, examining 21 compounds across the compositionality spectrum.
Our three experiments reveal that compound concepts are not stored as dedicated directions.
Instead, the model uses three complementary mechanisms: (1) contextual modification of head noun representations (cosine similarity 0.954), (2) compound-specific SAE features that capture fine-grained differences invisible to geometric measures (56.3\% unique features), and (3) massive next-token prediction priming that constructs compounds sequentially ($4{,}221\times$ for ``washing machine'').

The key takeaway is that LLMs solve the ``too many concepts, too few dimensions'' bottleneck for compound concepts by distributing representation across time steps rather than packing more features into the same vector space.
This is a distinct strategy from superposition: where superposition handles single-token features, sequential construction handles multi-token concepts.

Future work should extend this analysis to larger models where dedicated compound directions may emerge, test with richer natural contexts, apply causal interventions to establish whether compound-specific SAE features are functionally necessary, and explore whether the sequential construction pattern holds across languages with different compounding productivity.
